\documentclass{article}
\usepackage[T1]{fontenc}
\usepackage{lmodern}
\usepackage{graphicx}
\usepackage{amsmath,amssymb}
\usepackage[a4paper,margin=2cm]{geometry}
\usepackage[absolute,overlay]{textpos}
\setlength{\TPHorizModule}{1mm}
\setlength{\TPVertModule}{1mm}
\usepackage[indonesian]{babel}
\usepackage{hyperref}
\setlength{\emergencystretch}{2em}
% unit: Halaman 17

\begin{document}

% === Halaman 17 (llm) ===

\begin{itemize}
  \item Pembuktian
  
  \"kunci OTP (acak) + plainteks (tidak acak) = cipherteks (acak)\" \\
  sebagai berikut:\
  Misalkan:\n  \[ C = (P + K) \pmod{26} \]
  dengan $K$ dipilih secara uniform:\n  \[ P(K = k) = \frac{1}{26}, \quad k = 0, 1, \dots, 25. \]
  Untuk setiap nilai plainteks tertentu $P = p$, pemetaan
  \[ C = (P + K) \pmod{26} \]
  merupakan pemetaan satu-ke-satu dari semua kemungkinan $K$ ke semua kemungkinan $C$.\
  Akibatnya:\n  \[ P(C = c \mid P = p) = \frac{1}{26} \]
  untuk setiap $c$ (yaitu: $P(C = \text{'a'}) = P(C = \text{'b'}) = \dots = P(C = \text{'z'}) = \frac{1}{26}$)
\end{itemize}

\end{document}
